%%%%%%%%%%%%%%%%%%%%%%%%%%%%%%%%%%%%%%%%%%%%%%%%%%%%%%%%%%%%%%
%%%%%%%%%%%%%%%% MA-0820 Teoría de modelos %%%%%%%%%%%%%%%%%%
%%%%%%%%%%%%%%%%%%%%%%%%%%%%%%%%%%%%%%%%%%%%%%%%%%%%%%%%%%%%%%
\newpage
\subsection*{MA-0820 Teoría de modelos}
\addcontentsline{toc}{subsection}{MA-0820 Teoría de modelos}

\hypertarget{MA0820}{}
\begin{refsection}
\begin{table}[H]
\centering{
\begin{tabular}{|ll|}
\hline
\textbf{Año:} IV  & \textbf{Requisitos:} MA-0711 Lógica \\ 
\textbf{Ciclo:} Optativa  & \textbf{Correquisitos:} Ninguno \\ 
\textbf{Tipo de curso:} Teórico & \textbf{Horas presenciales por semana:} 5 \\ 
\textbf{Créditos:} 5 & \textbf{Horas de trabajo independiente por semana:} 10 \\ \hline
\end{tabular}
}
\end{table}

\subsubsection*{Descripción del curso}

Este curso optativo del área de Lógica Matemática introduce la Teoría de Modelos, rama que estudia la relación entre estructuras matemáticas (grupos, cuerpos, grafos, espacios vectoriales, entre otras) y los lenguajes de primer orden con los que pueden describirse. La lógica matemática analiza los sistemas formales y las estructuras que los interpretan; la interacción entre sintaxis y semántica permite comprender qué propiedades se derivan de un conjunto de axiomas.

Un concepto central es el de \emph{conjunto definible}: el conjunto de soluciones de una fórmula de primer orden en una estructura. Esta noción generaliza naturalmente los conjuntos algebraicos definidos por ecuaciones polinomiales y es fundamental en aplicaciones de la teoría de modelos al álgebra, la geometría y otras áreas.

En la primera parte se repasan los fundamentos de la lógica de primer orden (sintaxis, semántica, completitud y compacidad). Posteriormente se introducen inmersiones elementales, equivalencia elemental, tipos, saturación y eliminación de cuantificadores, ilustrados mediante cuerpos algebraicamente cerrados, espacios vectoriales, órdenes densos, grupos y grafos. Se presentan aplicaciones representativas al álgebra, mostrando cómo las herramientas lógicas permiten obtener resultados matemáticos profundos y unificadores.

\subsubsection*{Objetivos}

\textbf{General:} Comprender los principios básicos de la teoría de modelos y de sus aplicaciones a las estructuras algebraicas.

\textbf{Específicos:} Al finalizar el curso se espera que la persona estudiante sea capaz de:

\begin{enumerate}
\item Entender el concepto formal de lenguaje y fórmula de primer orden y demostrar resultados sobre la sintaxis de un lenguaje.
\item Entender el concepto formal de estructura de primer orden y su relación semántica con las fórmulas.
\item Aplicar el teorema de compacidad a casos concretos de estructuras naturales.
\item Determinar si teorías básicas, como la de cuerpos algebraicamente cerrados, tienen eliminación de cuantificadores.
\item Comprender aplicaciones de la teoría de modelos a la geometría algebraica, como el Nullstellensatz y el 17.º problema de Hilbert.
\item Elaborar un proyecto de investigación que aplique una o varias de las herramientas del curso.
\end{enumerate}

\subsubsection*{Contenidos}

\begin{enumerate}

\item \textbf{Fundamentos de lógica de primer orden (3 semanas):}
\begin{enumerate}
    \item Repaso: lenguajes, términos, fórmulas, estructuras, satisfacción y consecuencia lógica.
    \item Teoremas de completitud y compacidad; aplicaciones básicas.
\end{enumerate}

\item \textbf{Principios básicos de teoría de modelos (3 semanas):}
\begin{enumerate}
    \item Teorema de Tarski--Vaught; subestructuras y extensiones elementales.
    \item Teoremas de Löwenheim--Skolem; categoricidad y criterio de Vaught.
    \item Método back-and-forth para equivalencias elementales e isomorfismos.
\end{enumerate}

\item \textbf{Eliminación de cuantificadores (4 semanas):}
\begin{enumerate}
    \item Criterios de eliminación de cuantificadores; relación con completitud y decidibilidad.
    \item Ejemplos: cuerpos algebraicamente cerrados, cuerpos real cerrados, órdenes densos sin extremos.
\end{enumerate}

\item \textbf{Aplicaciones al álgebra y la geometría algebraica (3 semanas):}
\begin{enumerate}
    \item Nullstellensatz de Hilbert y compacidad.
    \item Teoría de modelos y el 17.º problema de Hilbert (Artin).
    \item Reformulación de problemas algebraicos en términos de existencia de modelos.
\end{enumerate}

\item \textbf{Tipos y compacidad (3 semanas):}
\begin{enumerate}
    \item Tipos completos y parciales; relación con conjuntos definibles.
    \item Teorema de omisión de tipos; aplicaciones a la existencia de modelos.
    \item Extensión y realización de estructuras.
\end{enumerate}

\end{enumerate}

\subsubsection*{Metodología}

\noindent \textbf{Lineamientos metodológicos:} La persona docente a cargo de este curso debe respetar las siguientes pautas:

\begin{enumerate}
    \item Desarrollar el curso mediante \textbf{clases magistrales presenciales}, con participación activa del estudiantado.
    \item Complementar con \textbf{sesiones de ejercicios} para consolidar demostraciones y razonamiento lógico abstracto.
    \item Incluir \textbf{talleres o proyectos} (individuales o grupales) orientados a aplicaciones o profundización temática.
    \item Promover evaluación formativa y sumativa; se sugiere incluir al menos dos exámenes parciales y un proyecto.
    \item Cuando las autoridades sanitarias o institucionales impongan restricciones, adaptar las actividades según normativa vigente.
\end{enumerate}

\textbf{Sugerencias metodológicas:} Adicionalmente, se tienen las siguientes recomendaciones para la persona docente a cargo de este curso:

\begin{enumerate}
    \item Utilizar \cite{Mar02} y \cite{Poi00} como referencias principales del curso.
    \item Complementar con \cite{HiL19}, \cite{EFT94} y \cite{CL2000} para repaso de lógica de primer orden.
    \item Documentar la oferta y coordinar con MA-0711 Lógica para coherencia de prerrequisitos.
\end{enumerate}

\nocite{HiL19}
\nocite{CL2000}
\nocite{EFT94}
\nocite{Mar02}
\nocite{Poi00}

\printbibliography[heading=subbibliography,section=\the\value{refsection}]
\end{refsection}
